\documentclass{resume} % Use the custom resume.cls style

\usepackage[left=0.4 in,top=0.4in,right=0.4 in,bottom=0.4in]{geometry} % Document margins
\usepackage{hyperref}
\newcommand{\tab}[1]{\hspace{.2667\textwidth}\rlap{#1}} 
\newcommand{\itab}[1]{\hspace{0em}\rlap{#1}}
\name{Neeharika Yadlapalli} % Your name
% You can merge both of these into a single line, if you do not have a website.
\address{+91 8247485839 $\cdot$ Guntur, India} 
\address{\href{mailto:neeharika.yadlapalli@gmail.com}{Gmail} $\cdot$ \href{https://www.linkedin.com/in/neeharika-yadlapalli-459474265/}{LinkedIn}}


\begin{document}

%----------------------------------------------------------------------------------------
%	OBJECTIVE
%----------------------------------------------------------------------------------------

\begin{rSection}{SUMMARY}

 I am a passionate AI/ML Engineer driven by an obsession for building autonomous systems that push the limits of what Generative AI can do. Specializing in \textbf{Agentic Workflows}, \textbf{LangGraph}, and \textbf{LLM Optimization}, I don't just implement APIs, I architect intelligent solutions that solve complex, unstructured problems, earned \textbf{Second Runner-Up at the Guinness World Record-certified Google Cloud Agentic AI Day 2025}, and co-authored a research paper at \textbf{INNOCOMP-2024}. Focused on building cost-efficient autonomous systems that move beyond standard API integrations to solve complex, unstructured problems at scale.

\end{rSection}

%----------------------------------------------------------------------------------------
%	WORK EXPERIENCE SECTION
%----------------------------------------------------------------------------------------

\begin{rSection}{EXPERIENCE}

\textbf{Associate Technical Specialist, AI/ML} \hfill Oct 2025 - Present\\
\textbf{Engineering Intern, AI/ML} \hfill March 2025 - Oct 2025\\
DigitalAPI.ai \hfill \textit{Bangalore, India}
 \begin{itemize}
    \itemsep -3pt {} 
      \item \textbf{Engineered ``Analytics on Demand'' Platform:} Designed a natural language-to-analytics agent using \textbf{Gemini 2.0} and \textbf{Model Context Protocol (MCP)}; built \textbf{5+ custom MCP tools} for autonomous cross-database querying (\textbf{MongoDB}) across \textbf{millions of data rows}, deriving complex reasoning and generating dynamic visual dashboards in \textbf{under 3 seconds} per query — eliminating \textbf{15+ hours/week} of manual reporting at \textbf{100\% retrieval accuracy}.
     \item \textbf{ASK HELIX:} Developed an intelligent chatbot using LLM and \textbf{Model Context Protocol} (MCP) framework achieving \textbf{90\%}query resolution accuracy. Automated support for API management platform,enabling users to query API documentation, integration guidelines, and platform features through natural language processing.The solution enhanced developer experience by streamlining access to API specifications.
     \item \textbf{API GPT:} Developed API GPT, a centralized platform that automatically generates MCP servers using Gemini LLM on Vertex AI with \textbf{95\%}uptime and processes queries across 100+ integrated APIs.Built end-to-end system converting API contracts into deployable Python MCP servers with sub-500ms query routing, achieving \textbf{92\%} accuracy in API recommendations. Implemented GCS storage integration, smart query routing, subscription-based access control, and secure authentication via ASH system.
     \item \textbf{MCP Code Generation Server with Dynamic Visualization:} Engineered a comprehensive MCP (Model Context Protocol)server that dynamically generates, executes, and displays Python code \textbf{90\%} code execution success rate and average 3.2-second generation time based on natural language user queries, featuring integrated data visualization and real-time graphical representation capabilities.The solution provides end-to-end code generation workflow from query interpretation to execution output.
 \end{itemize}


\end{rSection}

%----------------------------------------------------------------------------------------
%	PROJECTS SECTION
%----------------------------------------------------------------------------------------

\begin{rSection}{FEATURED PROJECTS}
\vspace{-1.25em}

\item \textbf{Juno — Voice-Native Financial Operating System} | \textit{Google Cloud Agentic AI Day 2025} \\
\textit{Tech Stack: Flutter, Vertex AI (Gemini 2.5 Flash Lite), Firebase, Firestore, Cloud Run, Cloud Speech-to-Text, MCP}
\begin{itemize}
    \item \textbf{Multi-Agent Financial Reasoning:} Built a multi-agent advisory system in which a \textbf{Coordinator MCP} orchestrates specialized Gemini-powered agents (\textit{Growth, Security, Context, Feedback}); integrated with the \textbf{Fi Money MCP Server} to securely stream a user's normalized assets, liabilities, cash-flow, and risk metrics into agent context.
    \item \textbf{Voice-First Multilingual UX:} Engineered a \textbf{Flutter} mobile app with Cloud Speech-to-Text + TTS, enabling plain-language queries in multiple Indian languages — converting natural questions like \textit{``Could I retire at 50 if my SIP grows 12\%?''} into Monte-Carlo simulations and Search-grounded financial projections.
    \item \textbf{Production Cloud Backend:} Deployed serverless MCP micro-agents on \textbf{Cloud Run} with \textbf{Firebase Auth} (Phone-OTP, passkey), \textbf{Firestore} for per-user agent memory and life-event store (e.g., \textit{``wedding in Feb 2026 → high liquidity need''}), and \textbf{Google Search Grounding} for real-time market and regulatory data.
    \item \href{https://drive.google.com/file/d/12SFmxkaDG8kcLFpLTg6fAuKgeOlFq4Y8/view?usp=sharing}{\textit{Demo video}}
\end{itemize}
\end{rSection}

\begin{rSection}{KEY PROJECTS}
\begin{itemize}
\itemsep -3pt {}
    \item \textbf{Machine Learning-Based Digital Twin for Predictive Modeling in Wind Turbines:} Machine Learning-Based Digital Twin for Predictive Modelling in Wind Turbines Developed a 5G-NG-RAN assisted cloud-based digital twin framework using Microsoft Azure to virtually monitor wind turbines and predict power generation. Implemented a Gradient Boosting Trees (GBT) regression model on SCADA data to forecast wind speed and energy output,achieving \textbf{85\%} accuracy.
    \item \textbf{Detection of fake and clone accounts in twitter using classification,distance measure algorithms:} Built a machine learning-based system to detect fake and clone profiles on Twitter using Support Vector Machine (SVM) classification, analyzing user metadata, tweet content, and behavioral features to identify malicious accounts in Online Social Networks (OSN). 
\end{itemize}
\end{rSection}
\begin{rSection}{ACHIEVEMENTS} 
\begin{itemize}
     \item \textbf{Second Runner-Up — Google Cloud Agentic AI Day 2025:} Awarded \textbf{INR 1 Lakh} prize at the \textbf{Guinness World Record-certified} 30-hour innovation sprint by Google Cloud $\times$ Hack2Skill — the \textit{largest agentic-AI hackathon ever held}, with \textbf{1,941 adjudicated developers across 700+ teams}. Built \textit{Juno} (see Featured Projects) for the ``Let AI Speak to Your Money'' challenge. \href{https://drive.google.com/drive/folders/1cvr_w2xa_Q_uR6WmOFlj13R-buiMEQ7H?usp=sharing}{\textit{View Submission}}
   
    \item \textbf{Awarded as Best Outgoing Student, Department of IT:} Awarded Best Outgoing Student among B.Tech IT batch of 2020–2024 at VNITSW, recognizing top merit in academics, co-curriculars, and overall campus performance.
   \href{https://drive.google.com/drive/folders/1XbIEOGK1TiIXEYllnyi6C7TYoNTJv1rY?usp=sharing} {\textit{ View here}}

    \item \textbf {GRE Score: 319/340} | Quant: 168/170 (95th Percentile) | Verbal: 151/170 (2023)

   \item Co-authored "Dynamic Sign Language detection system using Media Pipe Holistic and LSTM based Deep learning Model" presented at the \textbf{IEEE conference} held at Sonipat under Emerging Innovations and Advanced computing organized by Novel Research Foundation and B.M Institute of Engineering and Technology, India- May 2024
  \href{https://ieeexplore.ieee.org/document/10664073}{\textit{View here}}

\end{itemize}
\end{rSection}

%----------------------------------------------------------------------------------------
%	EDUCATION SECTION
%----------------------------------------------------------------------------------------

\begin{rSection}{Education}

{\bf Bachelors in Information Technology} \hfill {2020 - 2024} \\
VIGNAN’S NIRULA Institute of Technology and Science for Women\hfill {CGPA - 8.8} \\
Guntur, India


\end{rSection}

%----------------------------------------------------------------------------------------
% TECHINICAL SKILLS	
%----------------------------------------------------------------------------------------
\begin{rSection}{TECHNICAL SKILLS}

\textbf{Languages}: JAVA, Python, SQL \\
\textbf{AI/ML \& Agentic}: LangChain, LangGraph, MCP (Model Context Protocol), RAG, Vector Search, TensorFlow, PyTorch, NumPy, Pandas, Matplotlib \\
\textbf{Databases}: MongoDB, MYSQL, BigQuery
\textbf{DevOps \& Tools}: GCP (Cloud Run, Vertex AI, Firebase, VMs), Git, Docker, Google Cloud Platform (GCP), Atlassian Suite (JIRA, Confluence)

\end{rSection}

%----------------------------------------------------------------------------------------
% SOFT SKILLS
%----------------------------------------------------------------------------------------
\begin{rSection}{SOFT SKILLS}

\textbf{Core Strengths}: Adaptability, Resilience, Team Collaboration, Time Management, Creativity. \\
\textbf{Spoken Languages}: English, Telugu and Hindi. \\

\end{rSection}

\end{document}
